% !TeX root = Thesis.tex

\chapter{Background}
\label{chap:background}

This chapter details the existing tools used in constructing CSG models and gauging their accuracy. We describe the functions used for creating SDF primitives, combining SDF primitives, and measuring reconstruction accuracy.

\section{SDF~Primitives}
\label{subsec:sdf_primitives}

CSG-CRN uses three shape primitives with SDF functions derived from~\cite{Quilez2011}. Given a query point~$p$, each function returns a signed distance value. The SDF functions themselves do not apply transformations. Instead, translation and rotation must be applied to the query point to represent a transformation of the SDF primitive.

The simplest SDF function is the ellipsoid
\begin{align*}
	\SDF_{\ellipsoid}(p) &= \frac{k_0 (k_0 - 1)}{k_1} \\
	k_0 &= \left\lVert \frac{p}{s} \right\rVert_2, \qquad
	k_1 = \left\lVert \frac{p}{s^2} \right\rVert_2,
\end{align*}
where~$s$ is the radius of the ellipsoid along each axis and~$\frac{p}{s}$ is an element-wise division.

The cuboid and cylinder SDF functions include a roundness variable~$\rho$, which is subtracted from the final SDF value because level sets farther from the object surface are rounder. This has a side effect of growing the primitive, so we multiply~$\rho$ by the scale of the object to maintain the size of the primitive. We scale by the smallest half-extent for cuboids and scale by the height for cylinders. The function for a cuboid SDF with rounding is
\begin{align*}
	\SDF_{\cuboid}(p) &= \lVert \max(u, 0) \rVert_2 + \min(\max(u_x, u_y, u_z), 0) - r \\
	u &= |p| - (s - r), \qquad r = \rho \min(s_x, s_y, s_z),
\end{align*}
where~$s$ is the half-extent of the cuboid along each axis. The cylinder SDF function with rounding is
\begin{align*}
	\SDF_{\cylinder}(p) &= \lVert \max(v, 0) \rVert_2 + \min(\max(v_1, v_2), 0) - r \\
	v &= \left( \frac{k_0 (k_0 - 1)}{k_1} + r, \; |p_z| - (s_z - r) \right), \qquad r = \rho s_z \\
	k_0 &= \left\lVert \frac{(p_x, p_y)}{(s_x, s_y)} \right\rVert_2, \qquad
	k_1 = \left\lVert \frac{(p_x, p_y)}{(s_x^2, s_y^2)} \right\rVert_2,
\end{align*}
where~$s_x$ and~$s_y$ are the semi-axes of the cylinder base and~$s_z$ is the half-extent of the cylinder height.

\section{Boolean~Operations}
\label{subsec:boolean_operations}

CSG-CRN can predict a union operation~$\min(a,b)$ or a subtraction operation~$\max(a,-b)$. In addition, the command includes a blending factor~$\alpha > 0$ to smooth out the result of the operation. We achieve this by implementing the smooth minimum~$\smin$ and smooth maximum~$\smax$ from~\cite{Quilez2013}. The functions for~$\smin$ and~$\smax$ are
\begin{align*}
	\smin_\alpha(a, b) &= \min(a, b) - \frac{\alpha}{4} h^2, \qquad
	h = \frac{\max(\alpha - |a - b|, 0)}{\alpha}, \\
	\smax_\alpha(a, b) &= -\smin_\alpha(-a, -b).
\end{align*}

\section{Marching~Cubes}
\label{sec:marching_cubes}

The SDF CSG output of CSG-CRN can be exported and directly used by programs such as \texttt{MagicaCSG}~\cite{Ephtracy2022}. Other use cases, such as computing accuracy metrics, require a discretized mesh output. This can be achieved using the Marching Cubes algorithm~\cite{Lorensen1987}. The signed distance is computed for each vertex of a grid encompassing the output bounds. Using the occupancy value of each grid point, the algorithm can determine which cells the surface passes through and generate a mesh face for each of those cells.

\section{Accuracy~Metrics}
\label{sec:accuracy_metrics}

There are several standard accuracy metrics used to measure the quality of a reconstruction. This section describes the formulas for the Chamfer Distance, Earth Mover's Distance, and Intersection Over Union metrics.

\subsection{Chamfer~Distance}
\label{subsec:chamfer_distance}

Chamfer Distance (CD) is a measure of the average distance error between two surfaces. A uniformly distributed set of surface samples~$S_1$ and~$S_2$ are gathered for the two surfaces. The metric becomes more accurate the more surface samples are used. The one-direction CD is computed by finding the closest sample in~$S_2$ for each sample in~$S_1$ and measuring the distance. This paper uses the bi-directional CD, which doubles the process by also finding the closest sample in~$S_1$ for each sample in~$S_2$~\cite{Fan2017}. The bi-directional CD is defined as
\begin{equation*}
	d_{\CD}(S_1, S_2) = \sum_{p \in S_1} \min_{q \in S_2} \lVert p - q \rVert_2^2 + \sum_{q \in S_2} \min_{p \in S_1} \lVert p - q \rVert_2^2, \qquad S_1, S_2 \subset \mathbb{R}^3.
\end{equation*}

\subsection{Earth~Mover's~Distance}
\label{subsec:earth_movers_distance}

Earth Mover's Distance (EMD) approximates the total distance samples on one surface need to be moved to reach the second surface. Similar to CD, a uniform set of surface samples~$S_1$ and~$S_2$ are needed for the surfaces being compared. Where CD allows multiple samples in~$S_1$ to be mapped to the same sample in~$S_2$, EMD requires a one-to-one bijection~$\phi$~\cite{Fan2017, Achlioptas2018}. The equation for EMD is
\begin{equation*}
	d_{\EMD}(S_1, S_2) = \min_{\phi : S_1 \to S_2} \sum_{p \in S_1} \lVert p - \phi(p) \rVert_2, \qquad S_1, S_2 \subset \mathbb{R}^3.
\end{equation*}

\subsection{Intersection~Over~Union}
\label{subsec:intersection_over_union}

The Intersection Over Union (IoU) metric is the intersection of two volumes~$A$ and~$B$ divided by their union. We approximate the intersection and union volumes using a set of uniformly sampled points~$U$. The IoU score decreases for each sample where the occupancy values of~$A$ and~$B$ disagree, and increases for each sample where the occupancy values agree~\cite{Mescheder2019}. IoU is defined as
\begin{equation*}
	\IoU(A, B) = \frac{|A \cap B|}{|A \cup B|} \approx \frac{\sum_{i=1}^{N} \mathbf{1}[\SDF_A(U_i) \le 0 \wedge \SDF_B(U_i) \le 0]}{\sum_{i=1}^{N} \mathbf{1}[\SDF_A(U_i) \le 0 \vee \SDF_B(U_i) \le 0]},
\end{equation*}
where~$N$ is the number of sample points in~$U$.
